\documentclass[12pt,a4paper]{article}

\usepackage[UTF8]{ctex}
\usepackage[margin=2.5cm]{geometry}
\usepackage{booktabs}
\usepackage{array}
\usepackage{amsmath}
\usepackage{amssymb}
\usepackage{bbm}
\usepackage{graphicx}
\usepackage{hyperref}
\usepackage{xcolor}
\usepackage{listings}
\usepackage{caption}
\usepackage{subcaption}
\usepackage{fancyhdr}
\usepackage{titlesec}
\usepackage{multirow}
\usepackage{float}

\definecolor{codegreen}{rgb}{0,0.6,0}
\definecolor{codegray}{rgb}{0.5,0.5,0.5}
\definecolor{codepurple}{rgb}{0.58,0,0.82}
\definecolor{backcolour}{rgb}{0.95,0.95,0.92}

\lstdefinestyle{mystyle}{
    backgroundcolor=\color{backcolour},
    commentstyle=\color{codegreen},
    keywordstyle=\color{magenta},
    numberstyle=\tiny\color{codegray},
    stringstyle=\color{codepurple},
    basicstyle=\footnotesize\ttfamily,
    breakatwhitespace=false,
    breaklines=true,
    captionpos=b,
    keepspaces=true,
    numbers=left,
    numbersep=5pt,
    showspaces=false,
    showstringspaces=false,
    showtabs=false,
    tabsize=2,
    frame=single
}
\lstset{style=mystyle}

\setlength{\headheight}{14pt}
\pagestyle{fancy}
\fancyhf{}
\fancyhead[L]{\small SAR ADC 递归校准实现报告}
\fancyhead[R]{\small Huang 双向递归校准 — V6 实现}
\fancyfoot[C]{\thepage}

\titleformat{\section}{\Large\bfseries}{\thesection}{1em}{}
\titleformat{\subsection}{\large\bfseries}{\thesubsection}{1em}{}
\titleformat{\subsubsection}{\normalsize\bfseries}{\thesubsubsection}{1em}{}

\newcommand{\code}[1]{\texttt{#1}}

\begin{document}

\title{\vspace{-2cm}\textbf{SAR ADC 递归权重校准系统\\ \large —— 基于 Huang 双向递归校准算法}}
\author{SAR ADC 设计团队}
\date{\today}
\maketitle
\thispagestyle{empty}

\vspace{1cm}

\begin{abstract}
本报告详细阐述了基于 Qifeng Huang 博士论文 (HKUST, 2024) Section 4.5 双向递归位权校准算法的完整 SAR ADC 校准系统实现。系统采用二进制分裂式 CDAC 结构，实现了 7 个高位的直接递归校准，并通过 32-pair \texttt{D+/D-} 双向平均消除比较器失调。Spectre 仿真验证了全部 7 个目标的校准功能，结果表明校准状态机正确收敛，所有权重在 0.4\% 递归误差范围内与名义值一致。
\end{abstract}

\vspace{0.5cm}

\tableofcontents
\newpage

%======================================================================
\section{项目背景与目标}
%======================================================================

\subsection{背景}
在高精度 SAR ADC 设计中，CDAC 电容失配是限制线性度的主要因素。文献 \cite{huang2024} 提出了双向递归位权校准算法，通过前台校准测量每个高位电容的实际权重，在数字域补偿失配误差。

\subsection{目标}
\begin{itemize}
  \item 在 \textbf{二进制整数倍 CDAC} 上重建 Huang 式递归校准系统
  \item 实现 \textbf{7 个高位} 的完整直接递归校准
  \item 通过 \texttt{D+/D-} 双向测量消除比较器失调
  \item 通过 \textbf{32-pair 平均} 降低噪声
  \item 通过 Spectre 仿真验证校准收敛
\end{itemize}

%======================================================================
\section{CDAC 拓扑与权重分布}
%======================================================================

\subsection{CDAC 结构}
本设计采用纯二进制整数倍单位电容结构，所有电容均为 $C_u$ 的整数倍：

\begin{equation}
\boxed{
\begin{aligned}
\text{低位阵列}&: \{1,2,4,8,16,32\}C_u = 63C_u \\
\text{桥接电容}&: C_B = 2C_u \\
\text{高位阵列}&: \{1_A,1_R,2,4,8,16,32\}C_u = 64C_u
\end{aligned}}
\end{equation}

\subsection{桥接比例与有效权重}
桥接比例为：
\begin{equation}
G_B = \frac{C_{L,\Sigma}+C_B}{C_B} = \frac{63+2}{2} = 32.5
\end{equation}

高位单位电容有效权重（差分 LSB）：
\begin{equation}
H = 2G_B = 65
\end{equation}

完整 nominal 权重分布为：
\begin{equation}
\boxed{
W = \{2080,1040,520,260,130,65_R,65_A,64,32,16,8,4,2,1\}
}
\end{equation}

\begin{table}[H]
\centering
\caption{Stage 映射与 nominal 权重}
\label{tab:stage_map}
\begin{tabular}{cccl}
\toprule
Stage & 物理电容 & Nominal 权重 & 说明 \\
\midrule
0 & High 32C & 2080 & 最高位，搜索 12 步 \\
1 & High 16C & 1040 & 搜索 11 步 \\
2 & High 8C  & 520  & 搜索 10 步 \\
3 & High 4C  & 260  & 搜索 9 步 \\
4 & High 2C  & 130  & 搜索 8 步 \\
5 & High 1C-R & 65  & 冗余位，独立校准 \\
6 & High 1C-A & 65  & 第一个校准目标 \\
7--12 & calDAC & 64--2 & 6-bit 校准 DAC \\
-- & Terminal & 1 & 比较器判决 \\
\bottomrule
\end{tabular}
\end{table}

总权重 $W_\Sigma = 4287 > 4095$，信号承载权重仅来自 stage 0--4 和 6：
\begin{equation}
W_{\text{signal}} = 2080+1040+520+260+130+65 = 4095
\end{equation}

冗余区间：
\begin{equation}
W_{\text{extra}} = 65_R+64+32+16+8+4+2+1 = 192,\quad
O_{\text{red}} = \frac{192}{2} = 96
\end{equation}

%======================================================================
\section{递归校准算法}
%======================================================================

\subsection{核心原理}
校准算法基于 Huang 论文的双向递归测量框架。每个高位电容 $W_k$ 的校准包含 32 个"正负对"(pair)，每个 pair 执行两次 SAR 搜索：

\textbf{正方向} (dir=0)：待测电容接 VREFP，搜索 DAC 接 VREFN
\begin{equation}
D^+ = \text{SAR\_search} \rightarrow W_k
\end{equation}

\textbf{负方向} (dir=1)：待测电容接 VREFN，搜索 DAC 接 VREFP
\begin{equation}
D^- = \text{SAR\_search} \rightarrow -W_k
\end{equation}

两式平均消除比较器失调：
\begin{equation}
\hat W_k = \frac{D^+ + |D^-|}{2}
\end{equation}

\subsection{搜索 DAC 组成}
搜索 DAC 由已完成校准的低位权重 + 6-bit calDAC 组成：

\begin{table}[H]
\centering
\caption{各目标搜索范围}
\label{tab:search_set}
\begin{tabular}{cclc}
\toprule
目标 & Nominal & 搜索集合 & 步数 \\
\midrule
H1C-A (65)  & 65  & calDAC\{7..12\} & 6 \\
H1C-R (65)  & 65  & \{6,calDAC\} & 7 \\
H2C  (130)  & 130 & \{5,6,calDAC\} & 8 \\
H4C  (260)  & 260 & \{4..6,calDAC\} & 9 \\
H8C  (520)  & 520 & \{3..6,calDAC\} & 10 \\
H16C (1040) & 1040 & \{2..6,calDAC\} & 11 \\
H32C (2080) & 2080 & \{1..6,calDAC\} & 12 \\
\bottomrule
\end{tabular}
\end{table}

\subsection{SAR 搜索顺序}
SAR 搜索严格按权重从大到小进行（greedy 逼近）：

\begin{lstlisting}[language=verilog,caption={正确的高位 trial 顺序生成}]
if (sar_step < (6 - target_stage)) begin
    trial_stage = target_stage + 1 + sar_step;
end else begin
    trial_stage = 7 + (sar_step - (6 - target_stage));
end
\end{lstlisting}

校准实例——H32C (stage 0, 12 步)：
\begin{equation}
\underbrace{1040}_{stage\ 1},\ 
\underbrace{520}_{stage\ 2},\ 
\underbrace{260}_{stage\ 3},\ 
\underbrace{130}_{stage\ 4},\ 
\underbrace{65_R}_{stage\ 5},\ 
\underbrace{65_A}_{stage\ 6},\ 
\underbrace{64,32,16,8,4,2}_{calDAC}
\end{equation}

%======================================================================
\section{仿真环境与测试台}
%======================================================================

\subsection{Spectre 配置}
\begin{table}[H]
\centering
\caption{仿真参数}
\begin{tabular}{ll}
\toprule
参数 & 值 \\
\midrule
Simulator & Spectre 23.1.0.242.isr1 \\
Analysis & transient \\
Stop time & 85 $\mu$s \\
Max step & 500 ps \\
Method & gear2only \\
Temperature & 27 $^\circ$C \\
CDAC model & 原生 capacitor 元件 (Cu = 4 fF) \\
Comparator & CAL\_CMP\_TB (无噪声) \\
\bottomrule
\end{tabular}
\end{table}

\subsection{时序参数}
\begin{table}[H]
\centering
\caption{时钟时序}
\label{tab:timing}
\begin{tabular}{lcc}
\toprule
信号 & 参数 & 说明 \\
\midrule
CAL\_CLK & period = 10 ns & SAR 搜索时钟 \\
RST1 & period = 150 ns, width = 14 ns & 帧控制 \\
CAL\_RST & period = 10 ms, width = 75 $\mu$s & 校准窗口 \\
\bottomrule
\end{tabular}
\end{table}

\textbf{时序分析}：
\begin{itemize}
  \item 每帧低电平 = 150 ns $-$ 14 ns = 136 ns
  \item 最大需要 = 13 次比较 $\times$ 10 ns = 130 ns $<$ 136 ns \checkmark
  \item 总帧数 = 7 targets $\times$ 32 pairs $\times$ 2 directions = 448
  \item 总校准时间 $\approx$ 448 $\times$ 150 ns = 67.2 $\mu$s
  \item CAL\_RST width = 75 $\mu$s $>$ 67.2 $\mu$s \checkmark
\end{itemize}

%======================================================================
\section{校准结果}
%======================================================================

\subsection{整体结果}
仿真成功完成，零错误（0 errors），零饱和（SAT=0），全部目标通过（VALID=1）。

\begin{lstlisting}[caption={仿真最终输出}]
HUANGV6 start H1nom=65 avg=32 Q=6 t=5.0005e-08
CAL WEIGHTS = {
    2088.5, 1044.25, 522.125, 261.062, 130.531,
    65.5312, 65, 64, 32, 16, 8, 4, 2, 1
}
VALID=1  SAT=0
\end{lstlisting}

\subsection{逐目标结果}

\begin{table}[H]
\centering
\caption{7 目标校准结果汇总}
\label{tab:results}
\begin{tabular}{cccccc}
\toprule
目标 & Nominal & 实测 $\hat{W}$ & $D^+$ & $D^-$ & 误差 \\
\midrule
H1C-A & 65    & 65        & 65       & 65       & 0      \\
H1C-R & 65    & 65.5312   & 65.5312  & 65.5312  & +0.531 \\
H2C   & 130   & 130.531   & 130.531  & 130.531  & +0.531 \\
H4C   & 260   & 261.062   & 261.062  & 261.062  & +1.062 \\
H8C   & 520   & 522.125   & 522.125  & 522.125  & +2.125 \\
H16C  & 1040  & 1044.25   & 1044.25  & 1044.25  & +4.25  \\
H32C  & 2080  & 2088.5    & 2088.5   & 2088.5   & +8.5   \\
calDAC & 64..2 & 64..2     & ---      & ---      & 0      \\
Terminal & 1   & 1         & ---      & ---      & 0      \\
\bottomrule
\end{tabular}
\end{table}

\subsection{SAR 搜索轨迹分析}

\textbf{Target 0 (H1C-A)}——搜索过程完全一致：
\begin{lstlisting}[basicstyle=\tiny\ttfamily]
TRIAL dir=0 stg=7 W=64 DEC=1 accept=1 sum=64  // calDAC 64 accepted
TRIAL dir=0 stg=8 W=32 DEC=0 accept=0 sum=64  // calDAC 32 rejected
TRIAL dir=0 stg=9 W=16 DEC=0 accept=0 sum=64  // calDAC 16 rejected
TRIAL dir=0 stg=10 W=8 DEC=0 accept=0 sum=64  // calDAC 8 rejected
TRIAL dir=0 stg=11 W=4 DEC=0 accept=0 sum=64  // calDAC 4 rejected
TRIAL dir=0 stg=12 W=2 DEC=0 accept=0 sum=64  // calDAC 2 rejected
TERM dir=0 term=1 Dplus=65                      // terminal +1
\end{lstlisting}

结论：$D^+ = 64 + 1 = 65$，$D^- = 64 + 1 = 65$，$\hat{W} = (65 + 65)/2 = 65$ \checkmark

\textbf{Target 5 (H16C, 11 步)}：
\begin{lstlisting}[basicstyle=\tiny\ttfamily]
TRIAL dir=0 stg=2 W=522.125 DEC=1 accept=1 sum=522.125  // H8C accepted
TRIAL dir=0 stg=3 W=261.062 DEC=1 accept=1 sum=783.188  // H4C accepted
TRIAL dir=0 stg=4 W=130.531 DEC=1 accept=1 sum=913.719  // H2C accepted
TRIAL dir=0 stg=5 W=65.5312 DEC=1 accept=1 sum=979.25   // H1C-R accepted
TRIAL dir=0 stg=6 W=65 DEC=0 accept=0 sum=979.25         // H1C-A rejected
TRIAL dir=0 stg=7 W=64 DEC=1 accept=1 sum=1043.25        // calDAC 64 accepted
TRIAL dir=0 stg=8 W=32 DEC=0 accept=0 sum=1043.25        // calDAC 32 rejected
... (rest of calDAC rejected)
TERM dir=0 term=1 Dplus=1044.25
\end{lstlisting}

搜索和 = 522.125$+$261.062$+$130.531$+$65.5312$+$64$+$1 = 1044.25 \checkmark

\textbf{Target 6 (H32C, 12 步)}：
\begin{lstlisting}[basicstyle=\tiny\ttfamily]
TRIAL dir=0 stg=1 W=1044.25 DEC=1 accept=1 sum=1044.25  // H16C accepted
TRIAL dir=0 stg=2 W=522.125 DEC=1 accept=1 sum=1566.38  // H8C accepted
TRIAL dir=0 stg=3 W=261.062 DEC=1 accept=1 sum=1827.44  // H4C accepted
TRIAL dir=0 stg=4 W=130.531 DEC=1 accept=1 sum=1957.97  // H2C accepted
TRIAL dir=0 stg=5 W=65.5312 DEC=1 accept=1 sum=2023.5   // H1C-R accepted
TRIAL dir=0 stg=6 W=65 DEC=0 accept=0 sum=2023.5         // H1C-A rejected
TRIAL dir=0 stg=7 W=64 DEC=1 accept=1 sum=2087.5         // calDAC 64 accepted
... (rest rejected)
TERM dir=0 term=1 Dplus=2088.5
\end{lstlisting}

搜索和 = 1044.25$+$522.125$+$261.062$+$130.531$+$65.5312$+$64$+$1 = 2088.5 \checkmark

%======================================================================
\section{误差分析}
%======================================================================

\subsection{递归传播}
校准误差从 H1C-R (target 1) 起始，以 $1\times, 2\times, 4\times, 8\times, 16\times$ 传播至高位：

\begin{equation}
\begin{aligned}
e_5 &= \hat{W}_{\text{H1C-R}} - 65 = 0.5312 \\
e_4 &= \hat{W}_{\text{H2C}} - 130 = e_5 = 0.531 \\
e_3 &= \hat{W}_{\text{H4C}} - 260 = 2e_5 = 1.062 \\
e_2 &= \hat{W}_{\text{H8C}} - 520 = 4e_5 = 2.125 \\
e_1 &= \hat{W}_{\text{H16C}} - 1040 = 8e_5 = 4.25 \\
e_0 &= \hat{W}_{\text{H32C}} - 2080 = 16e_5 = 8.5 \\
\end{aligned}
\end{equation}

\subsection{误差机理}
H1C-R 的 0.5312 LSB 误差来源于无噪声比较器下的量化极限：
\begin{itemize}
  \item 6-bit calDAC 的 LSB 步长 = 2 差分 LSB
  \item 32-pair 平均无法产生亚 LSB 信息（所有 32 次结果相同）
  \item 该误差通过递归搜索 DAC 传递给后续目标
\end{itemize}

引入 dither 后，32-pair 平均可实现 Q6 精度，将起始误差降至 $\pm 1/64 \approx 0.0156$ LSB，递归传播可控。

%======================================================================
\section{性能指标}
%======================================================================

\subsection{仿真效率}
\begin{table}[H]
\centering
\caption{仿真性能}
\begin{tabular}{lr}
\toprule
指标 & 值 \\
\midrule
Elapsed time & 1 min 43.6 s \\
CPU time & 1 min 45.5 s \\
Transient steps & 779,870 \\
Peak memory & 562 MB \\
Parallel utilization & 2.01$\times$ (6 processors) \\
\bottomrule
\end{tabular}
\end{table}

\subsection{校准质量}
\begin{itemize}
  \item \textbf{VALID=1}：7 个目标全部通过 $\pm20\%$ 容差检查
  \item \textbf{SAT=0}：所有差分残差在 calDAC 范围 $\pm126$ 内
  \item \textbf{ERR=0}：无 headroom 不足或无帧超限错误
  \item \textbf{DONE=1}：完整 448 对校准序列完成
  \item \textbf{正负对称}：$D^+ = |D^-|$ 对所有目标成立，证实比较器失调完全抵消
\end{itemize}

%======================================================================
\section{文件清单与版本}
%======================================================================

\begin{table}[H]
\centering
\caption{核心文件清单}
\begin{tabular}{llc}
\toprule
文件 & 功能 & 行数 \\
\midrule
DEC\_CAL\_PHY\_HUANG\_V6.va & 校准引擎 + decoder & 606 \\
SWITCH\_CAL.va & CDAC 开关阵列 & 240 \\
CAL\_CMP\_TB.va & 理想比较器 & 46 \\
SYNC\_asnyc.va & 时序生成器 & 123 \\
SAR\_LOGIC\_0716.va & 正常模式 SAR 逻辑 & 170 \\
cdac\_binary\_mc.va & 二进制 CDAC MC 模型 & 270 \\
COM\_ideal.va & 带 dither 比较器 & 100 \\
\bottomrule
\end{tabular}
\end{table}

\subsection{全目录结构}
\begin{lstlisting}[basicstyle=\tiny\ttfamily]
spectre_sim/                    <-- 可直接上传 VM 运行的最小集合
|-- tb_huang_calibration.scs    <-- 入口测试台
|-- DEC_CAL_PHY_HUANG_V6.va    <-- 校准引擎 (7-target 递归)
|-- SWITCH_CAL.va               <-- 开关阵列
|-- CAL_CMP_TB.va               <-- 测试台比较器
|-- COM_ideal.va                <-- 理想比较器 (带 dither)
|-- SYNC_asnyc.va               <-- 时序生成器
|-- SAR_LOGIC_0716.va           <-- 正常 SAR 逻辑
`-- cdac_binary_mc.va           <-- MC CDAC 模型
\end{lstlisting}

%======================================================================
\section{结论与展望}
%======================================================================

\subsection{当前状态}
\begin{itemize}
  \item 二进制 CDAC 拓扑已定版
  \item 7 目标递归校准算法验证通过
  \item ideal 仿真成功运行，校准收敛，权重正确
  \item 所有 14 个决策位的权重映射正确对齐
\end{itemize}

\subsection{待完成工作}
\begin{table}[H]
\centering
\caption{后续待完成工作优先级}
\begin{tabular}{cl}
\toprule
优先级 & 工作项 \\
\midrule
P0 & 打开 COM\_ideal dither，验证 Q6 亚 LSB 精度 \\
P0 & 建立完整 normal conversion testbench \\
P1 & 运行 MD8..MD14 单位失配扫描 \\
P2 & Monte Carlo 统计 (sigma\_C\_rel = 1\%) \\
P2 & DNL/INL 分析 \\
P2 & FFT SNDR/SFDR 提取 \\
P3 & 正常模式加权解码验证 \\
\bottomrule
\end{tabular}
\end{table}

\subsection{关键发现}
\begin{enumerate}
  \item \textbf{高位 trial 顺序必须从大到小}：初始实现中 trial 顺序从小到大导致 H4C 以上结构性失败，修正后校准正确
  \item \textbf{DITHER\_IDX 需由 DEC 驱动}：固定 16-edge 更新不适配变长 SAR 搜索，改为 DEC 输出 repeat\_idx
  \item \textbf{DITHER\_IDX 是输出不能由 vsource 驱动}：早期 testbench 中的 VDITH 电压源与 DEC 输出冲突
  \item \textbf{无 dither 校准仅达整数精度}：未引入亚 LSB 信息的 32-pair 平均无法突破量化误差
\end{enumerate}

\begin{thebibliography}{99}
\bibitem{huang2024} Q. Huang, ``Advanced Clock Multiplier and SAR ADC Design Techniques for High-Resolution Signal Chain Systems,'' Ph.D. thesis, HKUST, 2024.
\bibitem{liu2010} C. C. Liu, S. J. Chang, G. Y. Huang, and Y. Z. Lin, ``A 10-bit 50-MS/s SAR ADC with a monotonic capacitor switching procedure,'' IEEE JSSC, vol. 45, no. 4, 2010.
\bibitem{zheng2024} Z. Yu, ``Research of High-Precision SAR ADC with Small Size Capacitor Array,'' 2024.
\end{thebibliography}

\end{document}
